% =====================================================================
%  THE LEDGER FRAMEWORK — Market Data Manifesto, Version 1.0
%  Governing principles for market data. Subordinate to the
%  Architectural Constitution (ledger_manifesto_v1_4.tex).
%  Drafting lead: GATHERAL. Rigor-check: FORMALIS. Review: KLEPPMANN,
%  TALEB. Certification: CONCORDIA.
%  Prior art (source material, not a constraint): the Market Data
%  Calibration Manifesto (calibration_manifesto.tex).
%  Build: pdflatex (twice). Standard TeX Live; lmodern and microtype
%  used when available and skipped otherwise.
% =====================================================================
\documentclass[11pt,a4paper]{article}

\IfFileExists{lmodern.sty}{%
  \usepackage[T1]{fontenc}\usepackage{lmodern}%
  \IfFileExists{microtype.sty}{\usepackage{microtype}}{}%
}{}
\usepackage[margin=2.5cm,marginparwidth=2.0cm,marginparsep=2mm]{geometry}
\usepackage{parskip}
\usepackage[hidelinks]{hyperref}

% --- article addressing: margin identifiers, mirroring the Constitution ---
\IfFileExists{xcolor.sty}{\usepackage{xcolor}}{}
\ifdefined\textcolor\else\def\textcolor#1#2{#2}\fi
\newcommand{\mlabel}[1]{\leavevmode\hypertarget{md-#1}{}%
  \marginpar{\raggedright\scriptsize\ttfamily\textcolor{gray}{#1}}}

\setcounter{secnumdepth}{0}
\providecommand{\tightlist}{\setlength{\itemsep}{2pt}\setlength{\parskip}{0pt}}

% traceability tag at the close of each article
\newcommand{\trace}[1]{\par\smallskip{\small\itshape (#1)}}

\begin{document}

\begin{center}
  {\LARGE\bfseries THE LEDGER FRAMEWORK}\\[0.9em]
  {\large\bfseries Market Data Manifesto}\\[0.6em]
  {\normalsize Governing Principles for Market Data ---
   Subordinate to the Architectural Constitution}\\[1.4em]
  Version 1.0 \quad---\quad drafted July 2026 \quad---\quad certification pending
\end{center}

\vspace{0.6em}
\begin{center}\rule{0.35\linewidth}{0.4pt}\end{center}
\vspace{0.6em}

\begin{center}\itshape
The record holds observations; every price beyond them is derived from them.
\end{center}

\vspace{0.4em}

\begin{center}\textbf{Abstract}\end{center}
\begin{quote}
The Ledger records observations, not prices. A market data fact enters
the record exactly once, as a moveless observation admitted through the
one door; every curve, surface, correlation, and calibration is a
deterministic function of recorded inputs, of two kinds the document
keeps apart. A \emph{projection} computes over the record and rebuilds
from it alone; a model's output \emph{re-enters as a recorded
observation} (C-14.9, C-14.15). Reading back any recorded value is
unconditional; re-deriving a model's output is what needs the model. This
manifesto states the principles that make
append-only, event-driven, projected market data the only sensible way
to hold market data inside the Ledger.
Each principle is an \emph{article}, derived from six of the
Constitution's eight commitments --- auditability, reproducibility, time
travel, correctness, order, and simulability --- and resting on the
clause of the Constitution that carries it. The document is rated
against the Constitution, never the other way round; where the two
conflict, the Constitution prevails.
\end{quote}

\vspace{0.4em}
\begin{center}\rule{0.35\linewidth}{0.4pt}\end{center}
\vspace{0.6em}

\section{1. Market Data in the Ledger's Terms}

A market data fact is an \textbf{observation}: a recorded statement,
attributed to a source, that a named observable --- a price, a rate, a
fixing, an index level --- took a value with respect to a world-moment.
The observation is a fact about what a source said; it is not the truth
about what a thing is worth. Any ``true'' price behind the observation is
a modelling posit that lives in a model, never a claim the record makes.

The boundary is sharp. The source of an observation sits outside the
trust boundary, and the Events Executor captures the arrival and records
it (C-5.4). An observation enters the immutable log only as a moveless
\emph{observation-recording transaction} --- moveless because it records
data and moves no balances --- admitted through the single Transaction
Executor and folded into a \emph{home}, its resting place in the record
from which the contracts that need it read (C-4.8). (Throughout, to
\emph{fold} the log is to replay it in order, accumulating state; the
current state is the initial state plus the fold of all observations.)
There is no second door for market data: the one canonical record and the
single writer extend, without exception, to every observation. A kind of
market data is declared before any observation of it crosses --- its
fields, the market data operators that adjust it across events (C-9.2),
and the grain of the cause-derived identifier --- an identifier computed
from what caused an arrival, so a redelivery of the same cause collides
with it --- that decides when two arrivals are the same fact, all
registered on the record first, so an arrival of an undeclared kind is
quarantined, never guessed (C-4.12).

\textbf{Vocabulary, mapped once.} This manifesto uses the Constitution's
fixed names and coins no synonym. Its prior art, the Market Data
Calibration Manifesto, speaks a different dialect, mapped onto the
Constitution's vocabulary here once and not carried further:
\begin{itemize}
\tightlist
\item an \emph{Oracle} is an external source; the framework has no
  ``oracle'' primitive and introduces none (Article MD-3);
\item an \emph{attestation} is an observation, recorded by an
  observation-recording transaction (C-4.8);
\item \emph{event time} is execution time; \emph{knowledge time} splits
  into monitor time and door time (C-2.7);
\item a \emph{posterior} over parameters is a derived object (a value
  computed from recorded inputs; MD-6): a projection where its recipe
  computes over the record, a re-entered observation where it runs a
  model (C-14.15);
\item a \emph{model configuration} is the recorded, versioned terms of a
  derived object (C-4.4, C-7.2).
\end{itemize}
The Constitution's own term \emph{the market data operator} names one
thing only --- the corporate-action transform of C-9.2 that carries
pre-event data into the post-event frame --- and is never used for a
source.

\section{2. The Articles}

Each article derives from one or more of the six commitments and rests
on a clause of the Constitution. The identifiers MD-\emph{n} are
append-only: never reused, never renumbered, so later documents may cite
them.

\mlabel{MD-1}\textbf{MD-1. The observation is the atom of market data.}
A market data fact is an observation, and an observation is recorded
exactly once. It enters the immutable log as a moveless
observation-recording transaction, admitted through the single writer,
and folded into a home. ``Exactly once'' is a mechanism, not a wish: each
arrival carries its cause-derived identifier, and an arrival identical
under it is \emph{absorbed} as a duplicate, not recorded again (C-2.7,
C-12.3, C-14.3). Absorption is not a silent drop: no new value is
recorded, but the fact of the redelivery --- its arrival provenance and
monitor time --- is retained, so a healthy absorb is distinguishable from
a feed that has gone quiet. Three arrivals that look alike must be told
apart, and the identifier's grain --- a registered property of the data
kind (\S1) --- is what tells them: the same print \emph{redelivered} is
absorbed; a vendor \emph{correction} is a distinct, later observation
that names the value it corrects (MD-10); a genuinely \emph{new} print at
the same instant is a second observation, recorded in its own right. A
grain too \emph{fine} records a mere redelivery twice and double-counts;
a grain too \emph{coarse} absorbs a genuine second print and loses it.

That the registered grain is correct is a trust assumption, named beside
TA-ARRIVAL. The double-count is loud --- the extra observation is on the
record to be found --- but the over-coarse loss is \emph{silent}: the
second print is treated as a duplicate and never recorded, and where a
feed delivers two distinct prints identical under every field it sends,
no grain can separate them. This is the one residual case of silent loss,
and the manifesto does not pretend a grain always exists to prevent it;
it is caught where TA-ARRIVAL's gaps are caught --- at the perimeter, by
reconciling arrival counts against recorded observations. Subject to
that, nothing about market data is written twice: the value a source
stated is its sole canonical copy, and every price, curve, and surface
that mentions it is computed from it, never stored beside it.
\trace{Auditability, Correctness; C-1.4, C-4.8, C-2.7, C-12.3, C-14.3.}

\mlabel{MD-2}\textbf{MD-2. Capture precedes judgement.} The record
admits an observation before anything decides whether to believe it. A
stale, fat-fingered, or off-consensus quote is well-formed: it enters
through the normal door as an observation-recording transaction (C-4.8),
and judgement is deferred because economic correctness is never gated at
admission --- it is detected after the fact by recomputation (C-13.2,
C-13.3). A quote that is not even processable --- a corrupted payload,
an unresolvable reference --- is still not lost: it is captured with the
provenance of its arrival and stands as a visible open item (C-4.12).
Either way the Ledger captures, then classifies; it never refuses to
record in order to reject. Whether an observation is used, down-weighted,
or ignored is a separate and later question, answered by a derived
object over the record, never by the door. A ``rejected'' observation is
one a calibration assigned no weight --- present, inspectable, and open
to a different classification by a later derivation --- not one the
record never held. Capture-then-classify is a \emph{principle}, and its
guarantee runs \emph{from arrival}: whether an implementation's pipe can
hold the full firehose at the open is a capacity question outside this
document (trust assumption TA-ARRIVAL) --- but the principle forbids
using capacity as a licence to drop silently. What the boundary cannot
admit is a named, visible residual; what never reaches the boundary is
the perimeter's concern, reconciled there, never a fact the record can
claim (C-4.12).
\trace{Auditability, Order, Correctness; C-4.8, C-13.2, C-13.3, C-4.12.}

\mlabel{MD-3}\textbf{MD-3. The source is untrusted; its delivery is a
trust assumption.} The origin of an observation --- an exchange, a
broker, a polling service, a fixing provider, an on-chain feed --- is
machinery beyond the system's trust boundary: an external source, not one
of the system's own machines (C-5.1, C-5.8), and the framework grants no
source the standing of truth. That an observation
reaches the boundary at all is named as a trust assumption and
reconciled against external authorities at the perimeter; it is never a
fact the fold may assume (C-4.12). The framework therefore has no
privileged ``oracle'' component and introduces none: a source is where
an observation came from, recorded as provenance, and nothing more.
\trace{Auditability, Correctness; C-5.1, C-5.4, C-5.8, C-4.12.}

\mlabel{MD-4}\textbf{MD-4. Every observation bears three times, and as-of
and as-at are two of them.} By the Constitution an event carries an
execution time --- the world-moment its source asserts it happened, the
time a court would enforce, contestable only in the world and corrected
only by a later event; a monitor time --- when the boundary observed it;
and a door time --- when the single writer admitted it (C-2.7). Bitemporality is
not a feature bolted onto market data; it is these three times, already
present. The \emph{as-of} coordinate of an observation --- the
world-moment its value concerns --- is its execution time; the
\emph{as-at} coordinate --- when the record came to hold it --- is its
door time; monitor time is the further provenance that splits an
observation's lateness into the world's delay and ours. Because both
coordinates are recorded, the system computes both honest answers at any
moment: what was known as of a past date, and what is known now about
that date (C-12.1).
\trace{Time travel, Order; C-2.7, C-12.1.}

\mlabel{MD-5}\textbf{MD-5. Order is meaning, and the authoritative fold
is execution-ordered.} Observations are events, and their sequence is
load-bearing: a derived object depends on which observations preceded
it --- a surface fitted before a shock is not the surface fitted after
it. Meaning lives in execution order. An observation that arrives late,
whose execution time precedes observations already folded, takes its
place among them; every \textbf{projection} downstream refolds to include
it (C-2.7), while a \textbf{re-entered observation} whose input window it
lands in cannot re-fold itself --- the ledger runs no model (C-14.9) ---
and is instead flagged stale for re-derivation (MD-8). Neither is silent:
the insertion is flagged, every changed or staled value is flagged, and
the difference is published as a named explain item --- a line in the
profit-and-loss explain attributing the change to its cause (C-12.6). The
as-known view --- what the book showed with only the observations it then
held --- survives unrewritten beside the corrected one.
\trace{Order, Time travel; C-2.7, C-12.6, C-14.9.}

\mlabel{MD-6}\textbf{MD-6. Every derived object is a deterministic
function of recorded inputs.} A curve, a surface, a correlation, a
calibration is a function of what the record holds --- the observations
it consumes, any derived objects it consumes, each pinned to its own
version and \emph{cut} (the exact as-at boundary that fixes which
observations were in force; MD-12), the recipe that defines it, and,
where the recipe draws on randomness, a recorded seed. Two kinds of
derived object must be told apart, and the Constitution names both.

Where the recipe computes \emph{over the record} --- a discount factor
read off recorded quotes, an adjusted value from the market data operator
--- the object is a \textbf{projection}: it stores nothing and rebuilds
from the record alone, by any party, today or at a past date (C-2.2,
C-3.2, C-14.11). Where the recipe \emph{runs a model} --- a filter, an
optimiser, a fitted surface --- the ledger does not run it (C-14.9); the
model is evaluated outside the fold and its output re-enters through the
one door as a new \textbf{observation} (C-14.15). That number is then
stored, like any observation, and two senses of getting it back must not
be confused: \emph{reading it back} is always possible from the record
alone, because the value is on the record; \emph{re-deriving} it ---
showing the number follows from its inputs --- additionally needs the
retained model and the numerical environment it ran in, a governance
matter outside this scope. So a re-entered observation is read back
unconditionally and re-derived only with the model.

Either kind carries, as recorded provenance, a \textbf{complete lineage}:
the resolved input cut --- every observation and every derived object it
consumed, each pinned to version and cut --- and, for a model output, the
model version. Lineage reaches through the whole chain to the
observations at its leaves, so what an object depends on is always
legible from the record. A derivation that \emph{fails} is not an absence
either, but which kind of fact it is follows the same split. A
\textbf{projection} that fails --- a bootstrap with too few quotes at
this cut --- is itself a projection: its failure is legible by
recomputation, never stored as a standing fact, and refreshes to success
if a late arrival fills the gap (MD-5). A \textbf{model run} that
fails --- a fit that does not converge --- re-enters as a recorded
observation carrying its diagnostics and lineage, and is itself subject
to staleness if its inputs are later corrected. A failed \emph{capture}
(MD-2) is stored for a different reason: it is an input-side arrival, not
an output-side derivation. In every case an absent object never passes
for one never attempted. The demand is on the inputs, not the engine:
record enough ---
observations, derived inputs, recipe, seed, and the lineage that ties
them --- that the object is a deterministic function of the record.
\trace{Reproducibility, Simulability, Correctness; C-2.2, C-3.2, C-14.9,
  C-14.11, C-14.15.}

\mlabel{MD-7}\textbf{MD-7. The prior is a term; the posterior is a
derived object.} A Bayesian calibration is not banned; it is placed. The
prior --- including any assumed dynamics for how a parameter moves
between observations --- is a declared, recorded term. The posterior is
a derived object of that prior-term, the recorded observations, and the
recorded seed (MD-6): a projection where it is a closed computation over
the record, a re-entered observation where a filter or sampler evaluates
a model. Either way it is \emph{not} stored state --- no ``current
posterior'' sits beside the record and updates in place; the maxim that
today's posterior is tomorrow's prior names a recomputation over the
log, not a mutable variable. The framework mandates no dynamics --- a
random walk, a martingale, a mean reversion are modelling choices a
recipe may declare, never axioms the framework asserts --- and whatever
is declared is recorded, so the posterior is reproducible on MD-6's
terms.
\trace{Correctness, Reproducibility; C-2.8, C-4.11, C-14.15.}

\mlabel{MD-8}\textbf{MD-8. A broken state cannot hide.} A recurring
pathology of calibrated systems is the stored parameter that has drifted
from what the information implies: the system knows something it has not
yet written into its numbers. The Ledger defeats it, but the mechanism
differs by kind of derived object, and honesty requires naming both. A
\textbf{projection} stores nothing --- recomputed from the record on
every read (C-4.11) --- so it cannot lag the record at all: the
drifted-parameter state is unrepresentable, and under a joint recipe an
on-demand read is the correct joint conditional expectation, so even two
correlated legs cannot hold a divergent pair. A \textbf{re-entered
observation} does store a number: a fitted surface, or a joint
vol-and-correlation fit for a spread, sits on the record because the
ledger runs no model (C-14.9). It cannot drift in place --- a new fit is
a new observation, forward, naming the inputs it consumed (MD-6), never
an edit of the old --- so the system never quietly overwrites what it
knew. But it can fall out of date when an input it consumed is later
corrected or superseded (MD-10, MD-5), and there the ledger does not
silently recompute it. Instead its staleness is a recorded fact: the
recorded lineage shows the input moved under it, and the stale fit stands
as a flagged open item awaiting re-derivation. So the joint spread
surface cannot silently disagree with a corrected leg --- not because it
re-fits itself, but because its disagreement is visible. So this
visibility is reached by default: a consumer reads a re-entered
observation through a projection --- ``the current fit for X'', which
selects the latest non-superseded observation and carries its staleness
flag --- never the raw stored value, so an ordinary read inherits the
check. What the
framework forbids everywhere is the \emph{silent} divergence: for a
projection it cannot arise; for a re-entered observation it cannot hide.
\trace{Correctness; C-1.5, C-4.11, C-11.5, C-14.9.}

\mlabel{MD-9}\textbf{MD-9. Constraints live in the derived object;
no-arbitrage is declared, not assumed.} A recorded observation carries
no promise of mutual consistency --- a stale quote and a fresh one may
cross, and two sources may disagree; the record holds both. No-arbitrage,
smoothness, and monotonicity are constraints a derived object imposes on
its own output, not properties the record must have. A derived object
that cannot meet the constraints it declares records that it cannot ---
a visible diagnostic --- rather than failing to exist or, worse, serving
an inadmissible object as though it were sound. No-arbitrage is not
negotiable for the object that claims it, but it is enforced by
detection, not by admission: an arbitrageable surface served as
arbitrage-free is an economic error, caught after the fact by
recomputation, not prevented at a door (C-13.2, C-13.3). The discipline
is that the claim and its test are both on the record. Which
no-arbitrage conditions are correct, and whether a fitted object hedges,
are matters of model choice and lie outside this scope (C-Scope.11): the
manifesto mandates that a derived object's constraints be recorded and
their satisfaction be a recorded diagnostic, not which constraints a
modeller should choose.
\trace{Auditability, Correctness; C-13.2, C-13.3, C-Scope.11.}

\mlabel{MD-10}\textbf{MD-10. Corrections repair forward; the original is
never overwritten.} A market data fact found wrong is not edited. The
discovery is a new recorded observation that names the wrong one, and the
discrepancy stands as a visible open item. What happens downstream splits
by kind: every \textbf{projection} over the corrected fact recomputes
automatically on the next read; a \textbf{re-entered observation} that
consumed the wrong fact cannot re-run itself (C-14.9), so its recorded
lineage marks it stale and it stands as a flagged open item until a fresh
fit supersedes it forward (MD-8). A bulk retraction --- a vendor voiding
a whole day --- is many corrections at once; the explain may summarise
them as one named item rather than flood the book, so the signal is not
lost in the flags. Where money already moved on the strength of the wrong
fact, it moves back only as an authorised compensating transaction
(C-12.4). The observed value stays on the record exactly as observed.
Adjustments a later event forces on earlier data --- a split rescaling a
price, a dividend rescaling a cash-per-share figure --- are computed at
each read by the market data operator (C-9.2, C-9.3), which is intrinsic
to the pairing of the kind of data and the kind of event, declared once
and never improvised. The original observation is preserved; the adjusted
value is derived; projections are recomputed from operator-adjusted
inputs.
\trace{Auditability, Time travel, Correctness; C-9.3, C-12.4, C-12.6, C-14.9.}

\mlabel{MD-11}\textbf{MD-11. Simulated market data is real market data
under a different seed.} The state of the market data universe is the
initial state plus the fold of all recorded observations (C-2.8). A
stress scenario, a backtest, a Monte-Carlo path is not a separate system
but the same machine folding a generated stream of observations from a
recorded seed; production is simply the one path whose observations
happened to be real. Because the seed is the single non-record input and
it too is recorded, every simulated path replays exactly, and a
calibration run over simulated observations is reproducible on the same
terms as one run over real ones. The seed suffices as the sole non-record
input only because every model output along the path is itself captured
as a re-entered observation (MD-6); a path that ran a model but recorded
only the seed could not replay without that model.
\trace{Simulability, Reproducibility; C-2.8, C-10.1.}

\mlabel{MD-12}\textbf{MD-12. Derivation composes.} A derived object may
consume another derived object --- a volatility calibration reads a
discount curve, a spread model reads two surfaces. When it does, the
dependency --- which object, at which version and cut --- is itself part
of the recorded recipe (MD-6), so the composite is again a deterministic
function of the record, and the split of MD-6 composes with it. Along a
chain of \textbf{projections}, reconstruction runs the whole chain from
the record, bit-for-bit. Where the chain passes through a \textbf{re-entered
observation}, that node enters as a \emph{leaf}: its recorded value is
read back from the record, and re-deriving it needs its retained model
and numerical environment (MD-6) --- so reconstruction of the composite
is bit-for-bit given those, and read-back always. Either way, because
each input is fixed at its own recorded cut, never read from a mutable
store, the ``mid-update input'' cannot arise (MD-8): a curve feeding a
calibration is consistent with that calibration by construction, not by
scheduling.
\trace{Reproducibility, Correctness; C-2.2, C-14.11.}

\section{3. The Life of an Observation}

One observation can walk the whole picture, and the articles are easier
to hold once it has. Follow AAPL's official close on a Tuesday, needed by
an equity volatility surface.

\begin{itemize}
\item \textbf{Recorded once.} The exchange, an external source, publishes
  AAPL's Tuesday close of 150.20; the Events Executor captures it and the
  Transaction Executor admits a moveless observation-recording
  transaction. The value 150.20 is now on the record, once; a redelivery
  of the same print is absorbed, not recorded again (MD-1, MD-3).
\item \textbf{Three times.} The observation carries its execution time
  (Tue 16:00, the world-moment the close concerns --- its as-of
  coordinate), its monitor time (when the boundary saw it), and its door
  time (Tue 18:30, when it was admitted --- its as-at coordinate) (MD-4).
\item \textbf{Read back, not drifting.} Tuesday's surface consumes
  150.20. If it is a projection it rebuilds from the record on read; if
  it is a fitted model it re-enters as a recorded observation pinned to
  the closes it consumed (MD-6). Either way its value is always readable
  back from the record; re-deriving a fitted surface needs its retained
  model. No surface is kept as a mutable copy that drifts (MD-8).
\item \textbf{Corrected forward.} On Thursday the exchange restates the
  Tuesday close to 150.50. This is a correction, not a late arrival: a
  new observation naming the old one, sharing Tuesday's execution time,
  separated only by door time (Thu) (MD-10). The original 150.20 stays on
  the record exactly as observed.
\item \textbf{Two honest answers, two numbers.} Now the daily question.
  \emph{Tuesday's surface as published on Tuesday} was built at the cut
  as-at Tue 18:30 --- it used 150.20, and that view survives unrewritten.
  \emph{Tuesday's surface as recomputed now} is built at the cut as-at
  today --- it uses the corrected 150.50. Same as-of date, two as-at
  cuts, two different surfaces; the cut is the coordinate that selects
  which (MD-4). A projection yields both by recomputation; a fitted
  surface yields the first as the observation recorded then and the
  second only once re-fitted, its staleness flagged until it is (MD-8).
\item \textbf{Late arrival.} Suppose instead a Tuesday-stamped trade
  print arrives only Thursday. Its execution time places it among
  Tuesday's data, so every projection over that window refolds to include
  it and publishes a named explain item, while a fitted surface for that
  window is flagged stale for re-derivation (MD-5). A value learned
  Thursday changes what Tuesday's book should have said.
\item \textbf{Simulated the same way.} Replace Tuesday's real close with
  a generated one from a recorded seed and the same machine runs
  unchanged; the surface reproduces on identical terms (MD-11).
\end{itemize}

\section{4. What This Manifesto Does Not Govern}

This manifesto governs how market data is recorded, ordered, corrected,
and projected. It does not govern:
\begin{itemize}
\item \textbf{price formation and price discovery} --- where a price
  comes from in the market is not a question the record answers;
\item \textbf{which model, prior, or no-arbitrage condition is correct}
  --- that is model choice, and its test is hedging performance, outside
  this document;
\item \textbf{model governance, validation, and the retention of
  models} --- a value reproduces from recorded prices unconditionally, a
  model's number only if the model is kept, and keeping it is
  governance;
\item \textbf{settlement} --- the record projects committed activity
  toward it but does not perform it.
\end{itemize}
These lie outside the Ledger boundary and are reconciled against it only
at its perimeter (C-Scope.11).

\section{5. Subordination}

This manifesto is subordinate to the Architectural Constitution and is
rated against it, never the other way round. Its vocabulary is the
Constitution's, one name per component, and it introduces no synonym for
any fixed term. Where a principle here would conflict with the
Constitution, the Constitution prevails, and the conflict is parked with
the exact text of the proposed amendment and the exact clause it would
replace --- never resolved by quiet narrowing. A departure from this
document, once genuine, is made by amending it explicitly first; an
article identifier is never reused and never renumbered.

\section{Conclusion}

Reconciliation of market data --- one desk's surface disagreeing with
another's, today's mark disagreeing with the mark reproduced tomorrow ---
has the same root cause as every other reconciliation: the same fact
stored more than once. The Ledger removes it here as everywhere. An
observation is recorded once; every curve, surface, and calibration is a
deterministic function of what the record holds; corrections repair
forward and originals survive; and the whole market data universe is the
initial state folded over its recorded observations. Operate market data any
other way inside the Ledger and you reintroduce the divergence the
architecture exists to remove. Operate it this way and there is nothing
left to reconcile.

\end{document}
